\documentclass[11pt]{article}
% commande d'exécution simple :
% pandoc RAPPORT.md --template=template_jpp12.tex -o RAPPORT.pdf --pdf-engine=xelatex
% pandoc RAPPORT.md --template=template_jpp12.tex --listings -o RAPPORT.pdf --pdf-engine=xelatex
% pandoc Rapport.md --template=template_jpp12.tex --listings --citeproc -o rapport.pdf --pdf-engine=xelatex
% pour avoir les emojis facilement, utiliser lualatex
% pandoc RAPPORT.md --template=template_jpp12.tex --listings -o RAPPORT.pdf --pdf-engine=lualatex
% ==================== ENCODAGE ET LANGUE ====================
\usepackage[T1]{fontenc}
\usepackage[utf8]{inputenc}
\usepackage[french]{babel}
\usepackage{fontspec}
\usepackage{xeCJK}
\frenchbsetup{StandardLists=true}
% ==================== POLICES PRINCIPALES ====================
\setmainfont{Libertinus Serif}
\setmonofont{Libertinus Mono}
\setCJKmainfont{Noto Serif CJK SC}
% ==================== EMOJIS COULEUR ====================
%\usepackage[unicode]{emoji}   % Active la saisie directe des emojis (✅, 👍, …)
%\setemojifont{Noto Color Emoji} % Police recommandée (présente sur Linux, Windows, macOS)
% ==================== MATHÉMATIQUES ====================
\usepackage{amsmath}
\usepackage{mathtools}
\usepackage{amssymb}
\usepackage{unicode-math}
\setmathfont{Libertinus Math}
\usepackage{dsfont}
% ==================== COMMANDES PERSONNALISÉES ====================
\usepackage{pifont}
\newcommand{\bra}[1]{\langle #1|}
\newcommand{\ket}[1]{|#1\rangle}
\newcommand{\braket}[2]{\langle #1|#2\rangle}
\newcommand{\R}{\mathbb{R}}
\newcommand{\C}{\mathbb{C}}
\newcommand{\N}{\mathbb{N}}
\newcommand{\Z}{\mathbb{Z}}
% Symboles de validation
\providecommand{\cmark}{\ding{51}}
\providecommand{\xmark}{\ding{55}}
\providecommand{\wmark}{\ding{73}}
\newcommand{\symSucces}{\textcolor{success}{\cmark}}
\newcommand{\symEchec}{\textcolor{failure}{\xmark}}
\newcommand{\symAttention}{\textcolor{warning}{\wmark}}
% ==================== TABLE DES MATIÈRES ====================
\usepackage{tocloft}
\setlength{\cftbeforesecskip}{2pt}
\setlength{\cftsecindent}{0pt}
\setlength{\cftsubsecindent}{10pt}
\setlength{\cftsubsubsecindent}{20pt}
\setcounter{tocdepth}{3}   % profondeur de la table des matières
% ==================== GÉOMÉTRIE ET MISE EN PAGE ====================
\usepackage{geometry}
\geometry{
margin=2.5cm,
headsep=4mm,
footskip=12pt
}
% Interlignes et espacements
\setlength{\baselineskip}{14pt}
\setlength{\parskip}{6pt}
\setlength{\parindent}{0pt}
\setlength{\parsep}{0pt}
\setlength{\itemsep}{2pt}
\setlength{\topsep}{2pt}
\setlength{\abovedisplayskip}{8pt}
\setlength{\belowdisplayskip}{8pt}
\setlength{\abovedisplayshortskip}{4pt}
\setlength{\belowdisplayshortskip}{4pt}
% ==================== TABLEAUX AMÉLIORÉS ====================
\usepackage{array}
\usepackage{booktabs}
\usepackage{longtable}
\usepackage{colortbl}
\usepackage{multirow}
\usepackage{float}
\usepackage{adjustbox}
\setlength{\LTpre}{6pt}
\setlength{\LTpost}{6pt}
\setlength{\extrarowheight}{2pt}
\renewcommand{\arraystretch}{1.2}
% ==================== LISTES ====================
\usepackage{enumitem}
\setlist{nosep, left=0pt}
% ==================== COULEURS ====================
\usepackage{xcolor}
\definecolor{success}{RGB}{0,150,0}
\definecolor{failure}{RGB}{200,0,0}
\definecolor{warning}{RGB}{255,140,0}
\definecolor{codebg}{RGB}{245,245,245}
% ==================== CODE ET LISTINGS ====================
\usepackage{listings}
\usepackage{framed}
%\usepackage{fancyvrb}   % optionnel
\definecolor{shadecolor}{rgb}{0.9,0.9,0.9}
% Les environnements Shaded et Highlighting seront définis par pandoc via 
\lstset{
basicstyle=\ttfamily\small,
breaklines=true,
frame=single,
backgroundcolor=\color{codebg},
numbers=left,
numberstyle=\tiny\color{gray},
keywordstyle=\color{blue},
commentstyle=\color{green!50!black},
stringstyle=\color{red},
showstringspaces=false,
tabsize=4,
language=Python
}
% ==================== BOÎTES ET ENCADRÉS ====================
\usepackage{tikz}
\usepackage[framemethod=tikz]{mdframed}
\usepackage{tcolorbox}
\tcbuselibrary{most}
\newtcolorbox{validationbox}[2][]{
colback=green!5!white,
colframe=green!75!black,
fonttitle=\bfseries,
title=#2,
#1
}
\newtcolorbox{errorbox}[2][]{
colback=red!5!white,
colframe=red!75!black,
fonttitle=\bfseries,
title=#2,
#1
}
\newtcolorbox{warningbox}[2][]{
colback=orange!5!white,
colframe=orange!75!black,
fonttitle=\bfseries,
title=#2,
#1
}
% ==================== FIGURES ET IMAGES ====================
\usepackage{graphicx}
\usepackage{adjustbox}
\usepackage{capt-of}
\usepackage{pdflscape}

% Paramètres par défaut
\setkeys{Gin}{keepaspectratio=true}

% Commande figure auto-cadrée
\newcommand{\autofigure}[3][]{
\begin{figure}[htbp]
    \centering
    \adjustbox{max width=\textwidth, max height=0.6\textheight, keepaspectratio, center}{
        \includegraphics{#2}
    }
    \captionof{figure}{#3}
    \label{fig:#2}
\end{figure}
}

% Commande figure paysage
\newcommand{\autolandscapefigure}[3][]{
\afterpage{
\begin{landscape}
\begin{figure}[p]
    \centering
    \adjustbox{max width=0.9\textwidth, max height=0.7\textheight, keepaspectratio, center}{
        \includegraphics{#2}
    }
    \captionof{figure}{#3}
    \label{fig:#2}
\end{figure}
\end{landscape}
}
}
% ==================== EN-TÊTE ET PIED DE PAGE ====================
\usepackage{fancyhdr}
% Style normal
\pagestyle{fancy}
\fancyhf{}
\fancyhead[C]{\textsc{Modèle Cl(6,6) — Rapport de Synthèse}}
\fancyfoot[C]{\thepage}
\renewcommand{\headrulewidth}{0.4pt}
\renewcommand{\footrulewidth}{0pt}
% Style "plain" = même numérotation que fancy
\fancypagestyle{plain}{
\fancyhf{}
\fancyhead[C]{\textsc{Modèle Cl(6,6) — Rapport de Synthèse}}
\fancyfoot[C]{\thepage}
\renewcommand{\headrulewidth}{0.4pt}
\renewcommand{\footrulewidth}{0pt}
}
% Première page sans numéro
\fancypagestyle{firstpage}{
\fancyhf{}
\fancyhead[C]{\textsc{Modèle Cl(6,6)}}
\fancyfoot[C]{}
\renewcommand{\headrulewidth}{0.4pt}
\renewcommand{\footrulewidth}{0pt}
}
% ==================== HYPERLIENS ====================
\usepackage{hyperref}
\hypersetup{
colorlinks=true,
linkcolor=blue,
urlcolor=blue,
citecolor=blue
}
% ==================== PROTECTION LIGNES ORPHELINES ====================
\clubpenalty=10000
\widowpenalty=10000
\displaywidowpenalty=10000
% ==================== DÉFINITIONS POUR PANDOC ====================
\providecommand{\tightlist}{%
\setlength{\itemsep}{2pt}\setlength{\parskip}{6pt}
}
\newcommand{\passthrough}[1]{#1}
% ==================== VARIABLES PANDOC ====================
\title{TRANSITIONS ANGULAIRES ENTRE PARTICULES}
\author{Bruno DE DOMINICIS}
\date{Mars 2026}
% ==================== REDÉFINITION DE \maketitle ====================
\makeatletter
\renewcommand{\maketitle}{
\begin{center}
\vspace*{1cm}
{\fontsize{18}{24}\selectfont \textbf{\@title} \par}
\vspace{1cm}
{\Large
\begin{tabular}[t]{c}
\@author
\end{tabular}\par}
\vspace{1cm}
{\Large \@date \par}
\end{center}
\vspace{1cm}
}
\makeatother
% ==================== CORPS DU DOCUMENT ====================
\begin{document}
\maketitle
\thispagestyle{firstpage}
% --- Mesure de la hauteur totale des abstracts ---
\newlength{\abstractheight}
\setbox0=\vbox{
\begin{center}
\begin{minipage}{0.85\textwidth}
\noindent \textbf{Résumé :} Nous présentons une reformulation des
transitions entre particules élémentaires (désintégration \(\beta\),
annihilation \(e^+e^-\), fusion nucléaire \(pp\), désintégration du
muon, oscillations de neutrinos, production de paires) comme des
réarrangements d'angles dans le réseau d'Ibozoo Uû, où chaque particule
est décrite par une pentade de l'algèbre de Clifford \(\text{Cl}(6,6)\).
Les générateurs \(i,j,k,I,J,K\) sont interprétés comme des angles. Nous
définissons un opérateur de transition \(T\) agissant sur l'espace de
Hilbert des pentades, décomposé en parties structure, feu, eau et mixte,
et nous établissons des règles de sélection (conservation des
générateurs, chiralité). Une formulation angulaire relie \(T\) à des
rotations dans l'espace à 6 dimensions. L'exemple de la désintégration
\(\beta\) est traité en détail. Ce formalisme offre une vision
géométrique unifiée des interactions, en accord avec les concepts
ummites.
\end{minipage}
\end{center}
\vspace{0.5cm}
\begin{center}
\begin{minipage}{0.85\textwidth}
\noindent \textbf{Abstract:} We present a reformulation of transitions
between elementary particles (\(\beta\) decay, \(e^+e^-\) annihilation,
\(pp\) nuclear fusion, muon decay, neutrino oscillations, pair
production) as angular rearrangements within the Ibozoo Uû network,
where each particle is described by a pentad of the Clifford algebra
\(\text{Cl}(6,6)\). The generators \(i,j,k,I,J,K\) are interpreted as
angles. We define a transition operator \(T\) acting on the Hilbert
space of pentads, decomposed into structure, fire, water and mixed
parts, and we establish selection rules (conservation of generators,
chirality). An angular formulation relates \(T\) to rotations in
6‑dimensional space. The example of \(\beta\) decay is treated in
detail. This formalism provides a unified geometric view of
interactions, in agreement with Ummite concepts.
\end{minipage}
\end{center}
\vspace{1cm}
}
\abstractheight=\ht0
\advance\abstractheight by \dp0
% Calcul de l'espace restant sur la page après le titre
\newlength{\remaining}
\remaining=\dimexpr\textheight - \pagetotal\relax
% Si la hauteur des abstracts dépasse l'espace restant, on agrandit la page
\ifdim\abstractheight > \remaining
\enlargethispage{\dimexpr\abstractheight - \remaining\relax}
\fi
\begin{center}
\begin{minipage}{0.85\textwidth}
\noindent \textbf{Résumé :} Nous présentons une reformulation des
transitions entre particules élémentaires (désintégration \(\beta\),
annihilation \(e^+e^-\), fusion nucléaire \(pp\), désintégration du
muon, oscillations de neutrinos, production de paires) comme des
réarrangements d'angles dans le réseau d'Ibozoo Uû, où chaque particule
est décrite par une pentade de l'algèbre de Clifford \(\text{Cl}(6,6)\).
Les générateurs \(i,j,k,I,J,K\) sont interprétés comme des angles. Nous
définissons un opérateur de transition \(T\) agissant sur l'espace de
Hilbert des pentades, décomposé en parties structure, feu, eau et mixte,
et nous établissons des règles de sélection (conservation des
générateurs, chiralité). Une formulation angulaire relie \(T\) à des
rotations dans l'espace à 6 dimensions. L'exemple de la désintégration
\(\beta\) est traité en détail. Ce formalisme offre une vision
géométrique unifiée des interactions, en accord avec les concepts
ummites.
\end{minipage}
\end{center}
\vspace{0.5cm}
\begin{center}
\begin{minipage}{0.85\textwidth}
\noindent \textbf{Abstract:} We present a reformulation of transitions
between elementary particles (\(\beta\) decay, \(e^+e^-\) annihilation,
\(pp\) nuclear fusion, muon decay, neutrino oscillations, pair
production) as angular rearrangements within the Ibozoo Uû network,
where each particle is described by a pentad of the Clifford algebra
\(\text{Cl}(6,6)\). The generators \(i,j,k,I,J,K\) are interpreted as
angles. We define a transition operator \(T\) acting on the Hilbert
space of pentads, decomposed into structure, fire, water and mixed
parts, and we establish selection rules (conservation of generators,
chirality). An angular formulation relates \(T\) to rotations in
6‑dimensional space. The example of \(\beta\) decay is treated in
detail. This formalism provides a unified geometric view of
interactions, in agreement with Ummite concepts.
\end{minipage}
\end{center}
\vspace{1cm}
% === AJOUT DE LA TABLE DES MATIÈRES ===
\tableofcontents
\newpage
% ======================================
\clearpage
\pagestyle{fancy}
\hypertarget{transitions-angulaires-entre-particules}{%
\section{TRANSITIONS ANGULAIRES ENTRE
PARTICULES}\label{transitions-angulaires-entre-particules}}

\hypertarget{remarque-sur-le-cadre-pruxe9-guxe9omuxe9trique}{%
\subsection{Remarque sur le cadre
pré-géométrique}\label{remarque-sur-le-cadre-pruxe9-guxe9omuxe9trique}}

Avant d'aborder les transitions entre particules, il est essentiel de
définir le cadre dans lequel elles se produisent. À l'instar des
approches qui recherchent des fondements algébriques ou géométriques
pour la physique fondamentale {[}1--4{]}, nous adoptons comme hypothèse
de travail l'idée que les particules élémentaires peuvent être décrites
comme des configurations stables au sein d'un substrat algébrique
pré-géométrique.

Dans cette approche, les degrés de liberté fondamentaux ne sont pas des
champs quantiques se propageant sur un fond spatio-temporel fixe, mais
plutôt des relations angulaires entre les générateurs d'une algèbre de
Clifford. Une analogie utile consiste à considérer le substrat comme un
\textbf{``faisceau d'axes orientés''} dont les orientations mutuelles
codent des informations physiques.

Il est essentiel de noter qu'un générateur isolé n'a aucune
signification physique directe ; c'est la structure relationnelle ---
\textbf{la configuration des angles entre les générateurs} --- qui
définit les grandeurs observables telles que la charge, la masse ou la
saveur.

Cette perspective est motivée par plusieurs considérations :

\begin{itemize}
\tightlist
\item
  Les algèbres de Clifford codent naturellement le spin, la chiralité et
  les structures de jauge {[}1, 2{]} ;
\item
  Les approches pré-géométriques suggèrent que l'espace-temps lui-même
  pourrait émerger de relations combinatoires ou algébriques plus
  fondamentales {[}3, 4{]} ;
\item
  Représenter les particules comme des ``motifs'' algébriques offre une
  voie pour unifier les interactions par des réarrangements géométriques
  plutôt que par des échanges de forces.
\end{itemize}

\textbf{Lien vers notre modèle des pentades} : Dans ce travail, nous
utilisons ce concept pour construire une représentation géométrique des
particules. Une particule n'est pas une entité immergée dans ce réseau ;
c'est une configuration locale et stable de ce réseau.

Notre formalisme dans Cl(6,6) et nos \textbf{pentades} constituent une
modélisation mathématique de cette idée : 1. Les six
\textbf{générateurs} \(i,j,k,I,J,K\) représentent les « axes »
fondamentaux de ce substrat algébrique dans notre cosmos. 2. Les
\textbf{angles} mentionnés tout au long du texte (par exemple,
\(\theta_1\theta_4\)) sont la traduction mathématique des relations
angulaires entre ces axes au sein du substrat. 3. Une \textbf{pentade},
cet ensemble de cinq éléments (trois bivecteurs de structure, un élément
axial « feu », un élément polaire « eau »), est la signature angulaire
complète qui définit l'état quantique et la nature d'une particule (sa
masse, sa charge, sa saveur) comme un motif spécifique et cohérent de
ces relations angulaires au sein du réseau.

Ainsi, lorsque nous décrivons une transition \(A \to B + C + \dots\)
comme un \textbf{``réarrangement des angles} au sein du \textbf{Réseau
angulaire fondamental''}, nous voulons dire que la configuration des
relations angulaires qui définissait la particule \(A\) se dissout et se
reconfigure pour former de nouvelles configurations stables, celles des
particules \(B\) et \(C\). La physique des particules devient ainsi une
\textbf{géométrie dynamique des angles}.

\begin{center}\rule{0.5\linewidth}{0.5pt}\end{center}

\hypertarget{principe-de-codification-des-pentades}{%
\subsection{Principe de codification des
pentades}\label{principe-de-codification-des-pentades}}

Dans notre modèle, chaque particule élémentaire est représentée par une
\textbf{pentade}, c'est‑à‑dire un ensemble ordonné de cinq éléments de
l'algèbre de Clifford \(\text{Cl}(6,6)\). Ces éléments sont interprétés
comme des \textbf{angles} dans un espace à six dimensions (générateurs
\(i,j,k,I,J,K\)). Ils se répartissent en trois rôles :

\begin{itemize}
\tightlist
\item
  \textbf{La structure} : trois bivecteurs (produits de deux
  générateurs) qui fixent l'identité profonde de la particule. Pour les
  nucléons, nous choisissons \(\{iI, jJ, kK\}\). Ce choix reflète la
  symétrie entre les trois dimensions spatiales (\(i,j,k\)) et les trois
  dimensions de charge (\(I,J,K\)). Il est commun au neutron et au
  proton car ils partagent la même composition en quarks de valence
  (deux quarks up et un quark down pour le proton\,; un up et deux down
  pour le neutron) mais avec des charges différentes.
\item
  \textbf{Le feu} : un élément de la forme \(i'v\), où \(i'\) est un
  générateur supplémentaire (souvent associé à la chiralité ou à
  l'interaction faible) et \(v\) un des six générateurs. Pour les
  nucléons, nous avons \(i'k\), qui représente la composante faible. Cet
  élément est identique pour le neutron et le proton car l'interaction
  faible agit de façon analogue sur les deux (ils peuvent se transformer
  l'un en l'autre).
\item
  \textbf{L'eau} : un élément de la forme \(1v\), où \(1\) est un
  élément scalaire (ou un générateur neutre) et \(v\) un générateur,
  représentant la \textbf{masse} ou la \textbf{substance} selon un axe
  privilégié. C'est cet élément qui porte la différence entre neutron et
  proton. La charge électrique est codée par l'orientation de cet axe :
  \(1i\) correspond à une masse alignée selon \(i\) (neutron, charge
  nulle), tandis que \(1j\) correspond à un alignement selon \(j\)
  (proton, charge positive). Cette interprétation est cohérente avec
  l'idée que la charge est une propriété directionnelle dans l'espace
  des angles.
\end{itemize}

\hypertarget{duxe9sintuxe9gration-ux3b2-n-p-e-ux3bdux2091}{%
\subsection{Désintégration β : n → p + e⁻ +
ν̄ₑ}\label{duxe9sintuxe9gration-ux3b2-n-p-e-ux3bdux2091}}

Ainsi, la désintégration β⁻ est vue comme un \textbf{réarrangement
angulaire} : le neutron change son angle de substance de \(i\) vers
\(j\), ce qui s'accompagne de l'émission de l'électron (qui possède
\(1j\)) et de l'antineutrino (qui possède \(1k\)). Les autres angles
(structure et feu) se recombinent pour former les pentades des
particules finales.

\hypertarget{pentades-initiale-et-finale}{%
\subsubsection{Pentades initiale et
finale}\label{pentades-initiale-et-finale}}

\textbf{État initial} (neutron) :

\[
P(n) = \{ iI,\ jJ,\ kK,\ i'k,\ 1i \}
\]

Angles associés (avec la correspondance \(i \leftrightarrow \theta_1\),
\(j \leftrightarrow \theta_2\), \(k \leftrightarrow \theta_3\),
\(I \leftrightarrow \theta_4\), \(J \leftrightarrow \theta_5\),
\(K \leftrightarrow \theta_6\)) : \[
\{\theta_1\theta_4,\ \theta_2\theta_5,\ \theta_3\theta_6,\ (\theta_1\theta_2\theta_3\theta_4\theta_5\theta_6)\cdot\theta_3,\ 1\cdot\theta_1\}
\]

\textbf{État final} (proton + électron + antineutrino) :

\[
P(p) = \{ iI,\ jJ,\ kK,\ i'k,\ 1j \}
\] \[
P(e^-) = \{ iI,\ iJ,\ iK,\ i'k,\ 1j \}
\] \[
P(\bar{\nu}_e) = \{ iK,\ iJ,\ iI,\ i'j,\ 1k \}
\]

(Nota : La représentation de l'antineutrino est ici prise identique à
celle du neutrino électronique utilisé dans la section sur les
oscillations ; une convention plus fine pourrait introduire des signes
pour distinguer particule et antiparticule, mais cela n'affecte pas la
discussion qualitative.)

\hypertarget{description-de-la-transition-angulaire}{%
\subsubsection{Description de la transition
angulaire}\label{description-de-la-transition-angulaire}}

\begin{itemize}
\tightlist
\item
  \textbf{Étape 1} : Le neutron change son angle de substance \(1i\) en
  \(1j\) (rotation de l'axe de masse de \(i\) vers \(j\)). C'est la
  transformation du quark down en quark up via l'interaction faible.
\item
  \textbf{Étape 2} : Les angles de structure du neutron
  \(\{iI, jJ, kK\}\) se réorganisent pour donner ceux du proton
  (inchangés) et ceux de l'électron (qui sont \(\{iI, iJ, iK\}\)).
  L'électron emporte une partie de la structure.
\item
  \textbf{Étape 3} : L'antineutrino est formé à partir des angles
  restants, avec une substance \(1k\) et un feu \(i'j\) (soit une
  permutation des générateurs).
\end{itemize}

Cette description géométrique remplace avantageusement l'image
habituelle de l'émission d'un boson \(W^-\) : la transition est un pur
réarrangement d'angles, sans introduction de particules virtuelles.

\begin{center}\rule{0.5\linewidth}{0.5pt}\end{center}

\hypertarget{annihilation-e-e-ux3b3ux3b3}{%
\subsection{Annihilation e⁺ e⁻ → γγ}\label{annihilation-e-e-ux3b3ux3b3}}

\hypertarget{principe-de-codification-des-pentades-pour-les-leptons-et-les-photons}{%
\subsubsection{Principe de codification des pentades pour les leptons et
les
photons}\label{principe-de-codification-des-pentades-pour-les-leptons-et-les-photons}}

\begin{itemize}
\tightlist
\item
  \textbf{La structure} : pour les leptons chargés, \(\{iI, iJ, iK\}\) ;
  pour le photon, idem mais sans feu ni eau.
\item
  \textbf{Le feu} : \(i'k\) pour l'électron, \(-i'k\) pour le positron.
\item
  \textbf{L'eau} : \(1j\) pour l'électron, \(-1j\) pour le positron.
\end{itemize}

Les photons sont représentés par \(\{iI, iJ, iK, 0, 0\}\).

\hypertarget{pentades-initiales-et-finales}{%
\subsubsection{Pentades initiales et
finales}\label{pentades-initiales-et-finales}}

\textbf{État initial} :

\[
P(e^-) = \{ iI,\ iJ,\ iK,\ i'k,\ 1j \}
\] \[
P(e^+) = \{ -iI,\ -iJ,\ -iK,\ -i'k,\ -1j \}
\]

\textbf{État final} (deux photons) :

\[
P(\gamma_1) = \{ iI,\ iJ,\ iK,\ 0,\ 0 \}
\] \[
P(\gamma_2) = \{ iI,\ iJ,\ iK,\ 0,\ 0 \}
\]

\hypertarget{description-de-la-transition-angulaire-1}{%
\subsubsection{Description de la transition
angulaire}\label{description-de-la-transition-angulaire-1}}

L'annihilation résulte d'une \textbf{cancellation des angles opposés} :
chaque élément du positron est l'opposé de celui de l'électron. Leur
somme donne zéro, mais l'énergie doit être conservée. Les six bivecteurs
(trois de chaque) se recombinent en deux ensembles de trois, formant
deux photons. Les éléments de feu et d'eau s'annulent réellement, car
les photons n'ont ni masse ni interaction faible.

En termes d'angles :

\[
\{iI,iJ,iK,i'k,1j\} + \{-iI,-iJ,-iK,-i'k,-1j\} \to \{iI,iJ,iK,0,0\} + \{iI,iJ,iK,0,0\}
\]

La conservation du nombre de générateurs est respectée : chaque photon
possède trois bivecteurs (six générateurs), tandis que la paire en avait
également six (avec signes opposés).

\begin{center}\rule{0.5\linewidth}{0.5pt}\end{center}

\hypertarget{fusion-nucluxe9aire-p-p-d-e-ux3bdux2091}{%
\subsection{Fusion nucléaire : p + p → d + e⁺ +
νₑ}\label{fusion-nucluxe9aire-p-p-d-e-ux3bdux2091}}

\hypertarget{pentades}{%
\subsubsection{Pentades}\label{pentades}}

\textbf{État initial} (deux protons) :

\[
P(p_1) = \{ iI,\ jJ,\ kK,\ i'k,\ 1j \}
\] \[
P(p_2) = \{ iI,\ jJ,\ kK,\ i'k,\ 1j \}
\]

\textbf{État final} (deutéron, positron, neutrino électronique) :

\[
P(d) = \{ iI,\ jJ,\ kK,\ i'k,\ 1j \} \quad (\text{avec liaison})
\] \[
P(e^+) = \{ -iI,\ -iJ,\ -iK,\ -i'k,\ -1j \}
\] \[
P(\nu_e) = \{ iK,\ iJ,\ iI,\ i'j,\ 1k \}
\]

\hypertarget{description}{%
\subsubsection{Description}\label{description}}

Un des protons subit une transformation β⁺ interne : son eau passe de
\(1j\) à \(1i\) (devenant ainsi un neutron), et il émet un positron
(configuration opposée) et un neutrino (structure permutée). Le neutron
ainsi formé se lie avec l'autre proton pour donner un deutéron. La
liaison est représentée par un couplage supplémentaire des angles, non
détaillé dans la pentade elle-même.

La conservation des générateurs est assurée par les règles de sélection
de l'interaction faible, qui autorisent les changements de saveur.

\begin{center}\rule{0.5\linewidth}{0.5pt}\end{center}

\hypertarget{duxe9sintuxe9gration-du-muon-ux3bc-e-ux3bdux2091-ux3bdux3bc}{%
\subsection{Désintégration du muon : μ⁻ → e⁻ + ν̄ₑ +
νμ}\label{duxe9sintuxe9gration-du-muon-ux3bc-e-ux3bdux2091-ux3bdux3bc}}

\hypertarget{pentades-1}{%
\subsubsection{Pentades}\label{pentades-1}}

\textbf{Muon} :

\[
P(\mu^-) = \{ jI,\ jJ,\ jK,\ i'i,\ 1k \}
\]

\textbf{Électron} :

\[
P(e^-) = \{ iI,\ iJ,\ iK,\ i'k,\ 1j \}
\]

\textbf{Antineutrino électronique} :

\[
P(\bar{\nu}_e) = \{ iK,\ iJ,\ iI,\ i'j,\ 1k \}
\]

\textbf{Neutrino muonique} :

\[
P(\nu_\mu) = \{ jK,\ jI,\ jJ,\ i'k,\ 1i \}
\]

(Remarque : la représentation des antineutrinos est ici calquée sur
celle des neutrinos, ce qui est une simplification. Une description plus
fine pourrait introduire des signes pour distinguer particule et
antiparticule, mais cela ne change pas l'interprétation géométrique de
la rotation d'axe.)

\hypertarget{description-1}{%
\subsubsection{Description}\label{description-1}}

Le muon a sa structure selon l'axe \(j\), l'électron selon \(i\). La
désintégration correspond à une \textbf{rotation de l'axe de saveur} de
\(j\) vers \(i\), accompagnée de l'émission de deux neutrinos qui
emportent les permutations. Les éléments de feu et d'eau se
redistribuent entre les trois produits finaux.

\begin{center}\rule{0.5\linewidth}{0.5pt}\end{center}

\hypertarget{oscillation-des-neutrinos-ux3bdux2091-ux3bdux3bc}{%
\subsection{Oscillation des neutrinos : νₑ →
νμ}\label{oscillation-des-neutrinos-ux3bdux2091-ux3bdux3bc}}

\hypertarget{pentades-2}{%
\subsubsection{Pentades}\label{pentades-2}}

\textbf{Neutrino électronique} :

\[
P(\nu_e) = \{ iK,\ iJ,\ iI,\ i'j,\ 1k \}
\]

\textbf{Neutrino muonique} :

\[
P(\nu_\mu) = \{ jK,\ jI,\ jJ,\ i'k,\ 1i \}
\]

\hypertarget{description-2}{%
\subsubsection{Description}\label{description-2}}

L'oscillation est une \textbf{rotation continue dans l'espace des
angles} : l'axe privilégié de la structure passe progressivement de
\(i\) à \(j\). L'opérateur de transition est \(\exp(i\alpha L_{ji})\),
avec \(\alpha\) proportionnel au temps et à \(\Delta m^2\). La
probabilité d'oscillation est \(\sin^2\alpha\), conforme à la formule
standard.

\begin{center}\rule{0.5\linewidth}{0.5pt}\end{center}

\hypertarget{production-de-paire-ux3b3-e-e}{%
\subsection{Production de paire : γ → e⁺ +
e⁻}\label{production-de-paire-ux3b3-e-e}}

\hypertarget{pentades-3}{%
\subsubsection{Pentades}\label{pentades-3}}

\textbf{Photon} (en présence d'un champ externe) :

\[
P(\gamma) = \{ iI,\ iJ,\ iK,\ 0,\ 0 \}
\]

\textbf{Paire électron-positron} :

\[
P(e^-) = \{ iI,\ iJ,\ iK,\ i'k,\ 1j \}
\] \[
P(e^+) = \{ -iI,\ -iJ,\ -iK,\ -i'k,\ -1j \}
\]

\hypertarget{description-3}{%
\subsubsection{Description}\label{description-3}}

Un photon de haute énergie interagit avec un champ externe (par exemple,
le champ coulombien d'un noyau). Ce champ fournit les orientations
manquantes (axes \(j\), \(k\), \(i'\)) qui permettent de \textbf{créer
la substance} : les éléments d'eau \(1j\) et \(-1j\) ainsi que les
éléments de feu \(i'k\) et \(-i'k\) apparaissent, tandis que la
structure du photon se scinde en deux ensembles opposés. C'est l'inverse
de l'annihilation.

\begin{center}\rule{0.5\linewidth}{0.5pt}\end{center}

\hypertarget{tableau-ruxe9capitulatif-des-transitions}{%
\subsection{Tableau récapitulatif des
transitions}\label{tableau-ruxe9capitulatif-des-transitions}}

\begin{longtable}[]{@{}
  >{\raggedright\arraybackslash}p{(\columnwidth - 4\tabcolsep) * \real{0.2857}}
  >{\raggedright\arraybackslash}p{(\columnwidth - 4\tabcolsep) * \real{0.1786}}
  >{\raggedright\arraybackslash}p{(\columnwidth - 4\tabcolsep) * \real{0.5357}}@{}}
\toprule
\begin{minipage}[b]{\linewidth}\raggedright
Transition
\end{minipage} & \begin{minipage}[b]{\linewidth}\raggedright
Type
\end{minipage} & \begin{minipage}[b]{\linewidth}\raggedright
Transformation angulaire
\end{minipage} \\
\midrule
\endhead
\(n \to p + e^- + \bar{\nu}_e\) & \(\beta^-\) & Rotation de l'angle de
substance (\(1i \to 1j\)) + émission \\
\(e^+e^- \to \gamma\gamma\) & Annihilation & Annulation d'angles
opposés \\
\(pp \to d e^+ \nu_e\) & Fusion & Entrelacement d'angles (avec
transformation d'un proton) \\
\(\mu^- \to e^- \bar{\nu}_e\nu_\mu\) & Désintégration & Rotation de
l'axe de saveur (\(j \to i\)) \\
\(\nu_e \rightleftarrows \nu_\mu\) & Oscillation & Rotation dans
l'espace des angles (dépendant du temps) \\
\(\gamma \to e^+e^-\) & Production de paire & Création de substance (eau
et feu) à partir d'angles, assistée par un champ externe \\
\bottomrule
\end{longtable}

\begin{center}\rule{0.5\linewidth}{0.5pt}\end{center}

\hypertarget{lois-de-conservation-angulaires}{%
\subsection{Lois de conservation
angulaires}\label{lois-de-conservation-angulaires}}

\begin{itemize}
\tightlist
\item
  \textbf{Conservation du nombre total de générateurs} : chaque élément
  (bivecteur, feu, eau) a un nombre fixe de générateurs ; la somme est
  invariante.
\item
  \textbf{Conservation du type de générateurs} : \(i,j,k,I,J,K\) (et
  \(i'\)) sont conservés individuellement, modulo transformations de
  jauge.
\item
  \textbf{Conservation de la chiralité} : liée à la présence de \(i'\)
  dans les éléments de feu.
\item
  \textbf{Conservation du moment angulaire total} : dans l'espace des
  angles, les rotations conservent la norme du vecteur angulaire.
\end{itemize}

\begin{center}\rule{0.5\linewidth}{0.5pt}\end{center}

\hypertarget{formalisation-mathuxe9matique}{%
\subsection{Formalisation
mathématique}\label{formalisation-mathuxe9matique}}

Soit \(\mathcal{H}\) l'espace de Hilbert des pentades (dimension 72 ou
144). Un opérateur de transition \(T\) agit sur les produits tensoriels
:

\[
\langle P_B \otimes P_C \otimes \cdots | T | P_A \rangle
\]

La probabilité de transition est donnée par l'élément de matrice au
carré.

\begin{center}\rule{0.5\linewidth}{0.5pt}\end{center}

\hypertarget{opuxe9rateur-de-transition-pentadique}{%
\section{OPÉRATEUR DE TRANSITION
PENTADIQUE}\label{opuxe9rateur-de-transition-pentadique}}

\hypertarget{duxe9composition}{%
\subsection{Décomposition}\label{duxe9composition}}

\[
T = T_{\text{structure}} + T_{\text{feu}} + T_{\text{eau}} + T_{\text{mixte}}
\]

\begin{itemize}
\tightlist
\item
  \(T_{\text{structure}}\) agit sur les triplets de bivecteurs.
\item
  \(T_{\text{feu}}\) sur les éléments \(i'v\).
\item
  \(T_{\text{eau}}\) sur les éléments \(1v\).
\item
  \(T_{\text{mixte}}\) couple les différents types.
\end{itemize}

\hypertarget{expression-compacte}{%
\subsection{Expression compacte}\label{expression-compacte}}

Dans \(\text{Cl}(6,0)\), avec une base \(\{E_a\}\) :

\[
T = \sum_{a,b} g_{ab} (E_a \otimes E_b^*)
\]

Pour une pentade \(P = \{B_1,B_2,B_3,F,S\}\) :

\[
|P\rangle = |B_1\rangle \wedge |B_2\rangle \wedge |B_3\rangle \wedge |F\rangle \wedge |S\rangle
\]

\[
T|P\rangle = \sum_{P'} T_{P'P} |P'\rangle,
\]

\[
T_{P'P} = \sum_{\text{cycles}} \prod_{k=1}^5 \langle e'_k | T_{\text{local}} | e_k \rangle
\]

\hypertarget{ruxe8gles-de-suxe9lection}{%
\subsection{Règles de sélection}\label{ruxe8gles-de-suxe9lection}}

\begin{itemize}
\tightlist
\item
  Conservation du nombre et du type de générateurs.
\item
  Conservation de la chiralité.
\item
  Diagrammes de Feynman pentadiques.
\end{itemize}

\hypertarget{exemple-duxe9sintuxe9gration-ux3b2}{%
\subsection{Exemple : désintégration
β}\label{exemple-duxe9sintuxe9gration-ux3b2}}

\[
\mathcal{M} = \langle P(p) \otimes P(e^-) \otimes P(\bar{\nu}_e) | T | P(n) \rangle = G_F \, J_{\text{had}} \cdot J_{\text{lep}}
\]

avec

\[
J_{\text{had}} = \langle P(p) | T_{\text{struct}} \otimes T_{\text{feu}} | P(n) \rangle,
\]

\[
J_{\text{lep}} = \langle P(e^-) \otimes P(\bar{\nu}_e) | T_{\text{feu}} \otimes T_{\text{eau}} |0\rangle.
\]

\hypertarget{quantification-canonique}{%
\subsection{Quantification canonique}\label{quantification-canonique}}

Opérateurs \(a_P\), \(a_P^\dagger\) avec
\([a_P, a_Q^\dagger]_{\mp} = \delta_{PQ}\). L'hamiltonien d'interaction
est une somme de produits.

\hypertarget{formulation-angulaire}{%
\subsection{Formulation angulaire}\label{formulation-angulaire}}

Les angles \(\theta_\mu\) sont les générateurs. \(T\) devient un
opérateur différentiel :

\[
T = \exp\!\left( i \sum_{i,j} \omega_{ij} L_{ij} \right), \quad L_{ij} = -i(\theta_i \partial_{\theta_j} - \theta_j \partial_{\theta_i})
\]

Pour la désintégration \(\mu \to e\), la probabilité est
\(\sin^2\alpha\), avec \(\alpha \propto t \Delta m^2/E\).

\begin{center}\rule{0.5\linewidth}{0.5pt}\end{center}

\hypertarget{vers-une-thuxe9orie-compluxe8te}{%
\subsection{Vers une théorie
complète}\label{vers-une-thuxe9orie-compluxe8te}}

Il reste à : - Déterminer les constantes de couplage à partir des
données. - Dériver les règles de Feynman pentadiques. - Calculer les
sections efficaces. - Prédire de nouveaux processus.

\hypertarget{conclusion}{%
\subsection{Conclusion}\label{conclusion}}

L'opérateur \(T\), défini par une somme sur les cycles d'éléments de
matrice locaux, fournit un cadre mathématique cohérent pour décrire
toutes les transitions entre particules comme des réarrangements
d'angles. Ce modèle géométrique, inspiré des concepts ummites, unifie
les interactions et ouvre la voie à une compréhension plus profonde de
la physique des particules.
\end{document}
